% LTeX: language=en-US
% !TEX root = main.tex
% LTeX: enabled=false
% IDE used for development: VS CODE with Latex Workshop
% TeX Live 2026
% Compile with LuaLatex

\documentclass[12pt, a4paper]{article}

% region — GENERAL PACKAGES o7
\usepackage{hyperref}       % link handling and pdf metadata
\usepackage{hyperxmp}       % link handling and pdf metadata
\usepackage{xurl}           % url handling
\usepackage{lua-ul}         % underlines, highlights, etc
\usepackage{luacolor}       % color overall
\usepackage{rotating}       % used for rotating :O (sidewaystable)
\usepackage{enumitem}       % used for custom enumeration (\begin{enumerate}[label=Question \arabic*:])
\usepackage{multicol}       % used for multicolumns
% endregion

% region — LANUGUAGE and locale settings 
\usepackage[main=english]{babel}   % language (naustrian)
\usepackage{csquotes}           % locale quotes (use \enquote{})
% endregion

% region — TITLE, METADATA and URL-Handling (hypersetup)
\title{
  Qualifying Spatial Relations using Jev: \\[0.5em]
  \large{What role could artificial system-one intelligence play in GeoAI?}\\[-3.5em]
}
\date{}
\author{}

\hypersetup{
    pdftitle={Qualifying Spatial Relations using Jev:},
    pdfauthor={Jakob Roland},
    pdflang={en},
    colorlinks=true,
    linkcolor=black,
    urlcolor=cyan,
    pdfborder={0 0 0}
}
\urlstyle{same} 
% endregion

% region — TYPESETTING (font, margins, formatting)
\usepackage{parskip}
\usepackage[hmargin=3cm, vmargin=1.5cm]{geometry}
\linespread{1.15}
\setlength{\headheight}{15pt}
% \setlength{\parindent}{1em} % to regain paragraph indent

% Font
\usepackage{fontspec}
\setmainfont{Noto Sans}
\newfontfamily\symbolsfont{DejaVu Sans}
\setmonofont{Source Code Pro}
% endregion

% region — OTHER FORM ELEMENTS (boxes, quotes, highlights)

% region — Boxes or Quotes
\usepackage[most]{tcolorbox}

% Intended for quotes standing out from the continous text 'simpleQuote' convinces with a minimalist design featuring only a introductory vertical line on the left
\newtcolorbox{simpleQuote}{
  enhanced jigsaw,
  colback=white,
  colframe=black!70,           
  arc=0mm,                 
  boxrule=0pt,    
  leftrule=3pt,                     
  top=8pt, bottom=8pt,     
  left=15pt, right=15pt,   
  width=\textwidth,         
  boxsep=0pt,
}
% endregion

% endregion

% region — REFERENCES (Zotero)
% \usepackage[style=apa]{biblatex}
\usepackage[style=numeric-comp]{biblatex}
% \usepackage[style=ieee]{biblatex}
\addbibresource{/home/jakob/ucloud/Zotero/MyLibrary_BibLatex.bib}
% endregion

% region — FLOATS (tables, images, code samples)
\usepackage{float}        % float placement, eg. [H] for here

% tables
\usepackage{tabularx}
\usepackage{array}
\usepackage{booktabs}
\usepackage{colortbl}
\newcolumntype{C}{>{\centering\arraybackslash}X}

% images
\usepackage{graphicx}
\graphicspath{ {./Figures/} }

% code samples (listings)
\usepackage{minted}
\renewcommand{\listingscaption}{Code Sample}
% endregion

% region — CUSTOM COMMANDS
% Emphasisized comments (Gemini chose the name)
\newcommand{\meth}[1]{\highLight[yellow]{\emph{#1}}}

% Context headings
\newcommand{\ctxH}[1]{
  {\hspace*{1cm}\textcolor[RGB]{180, 180, 180}{\textit{#1}}}
}
% endregion

% Latex Cheatsheet
\begin{comment}
Relative Size Commands
\tiny
\scriptsize
\footnotesize
\small
\normalsize
\large
\Large
\LARGE
\Huge
\end{comment}



\begin{document}

\maketitle
\textit{Disposition for a master thesis by Jakob Roland}\\
\textit{Supervised by Prof. Dr. Krzysztof Janowicz}\\
\textit{University of Vienna, 2026}

\section*{Introduction}
Since the release of the first chatbot powered by a large language model (LLM) in late 2022, it has become clear that artificial intelligence (AI) has left the realm of science fiction to enter our daily lifes. Since then almost all spheres of society have been impacted in some way by AI, Spatial Data Science being no exception. There have been many works testing LLMs in a variety of spatial applications [EXAMPLES] and some discussing their caveats, challenges, and shortcomings [EXAMPLES]. This thesis sees itself in both categories, recognizing concerns raised over the sustainability of GeoAI \cite{janowiczPhilosophicalFoundationsGeoAI2023,shiGeographyAISustainability2025} and evaluating a new variant of LLM for a spatial task. This thesis tests the spatial understanding of the system-one intelligence model Jev, which was released by the company Typesafe AI in September 2026. To this end, the experiment on spatial relation qualification done by Yuhan Ji (et al, 2025 \cite{jiFoundationModelsGeospatial2025}) is reproduced using Jev instead of ChatGPT-4. In contrast to ChatGPT-4 Jev is quicker, cheaper, and outputs in a fixed format, which facilitates downstream processing of the results. The question remains if the accuracy of Jev's answers is comparable to GPT-4, which this thesis will attempt to find out.

\end{document}
